\chapter{Text, tokenization, and sentiment systems}
\label{ch:text-sentiment}

\begin{learningobjectives}
After this chapter, you should be able to:
\begin{itemize}
  \item formulate sentiment analysis as a supervised prediction problem while
  stating what the label does and does not mean;
  \item distinguish text normalization, tokenization, vocabulary construction,
  integerization, padding, masking, and out-of-vocabulary handling;
  \item prevent preprocessing leakage by fitting a vocabulary on training data
  alone;
  \item explain embeddings, pooled text representations, binary logits, and
  cross-entropy mathematically;
  \item evaluate thresholded decisions with error-specific metrics and direct
  inspection of held-out failures; and
  \item specify the artifacts and monitoring signals required for a reliable
  text-classification system.
\end{itemize}
\end{learningobjectives}

\begin{chapterconnection}
The preceding chapters built nonlinear predictors from affine maps,
activations, objectives, and backpropagation. Text presents a new difficulty:
an observation is a variable-length sequence of symbols rather than a fixed
vector of measured real numbers. TensorFlow can differentiate and execute the
resulting model, but first we must define how language becomes tensors without
leaking held-out information or erasing important structure.
\end{chapterconnection}

\section{A scientific question before an architecture}

Consider a movie review $x$ and a binary label $y$, where $y=1$ denotes a
positive rating and $y=0$ denotes a negative rating. We seek a function whose
output supports the conditional probability claim
\[
  p_\theta(y=1\mid x).
\]
This is a supervised binary-classification problem. It is not a measurement of
an author's complete emotional state and it is not a universal definition of
positive language. The label inherits the data creator's rating threshold,
sampling decisions, time period, platform, and population.

The Stanford Large Movie Review Dataset is an instructive case. It contains
25,000 labeled reviews for training and 25,000 for testing, with balanced
positive and negative classes. Reviews with intermediate ratings are excluded.
That design makes a clean educational benchmark, but it also removes much of
the ambiguity that a deployed sentiment system would encounter.

\begin{definitionbox}{prediction contract for sentiment}
A sentiment prediction contract names the unit of analysis, label-construction
rule, represented population, information available at prediction time,
output semantics, decision threshold, error costs, and conditions under which
the model should abstain or defer to a person.
\end{definitionbox}

The first baseline need not be a neural network. A bag-of-words or TF--IDF
representation followed by logistic regression is fast, auditable, and often
competitive. A neural model should earn its added complexity through held-out
evidence or through a requirement---such as reusable learned representations
or sequence-sensitive prediction---that the baseline cannot satisfy.

\section{Data acquisition is part of the experiment}

The notebook obtains the archive from Stanford through Keras's cache-aware
\code{get\_file} utility. On the first student run, the utility downloads and
extracts the publisher's archive. Later runs reuse the local cache. An explicit
environment variable may instead identify an existing extracted
\code{aclImdb} directory. These are acquisition mechanisms, not alternative
datasets: both must resolve to the publisher's directory structure.

The path alone is insufficient provenance. An experiment record should retain
the publisher URL, retrieval date, archive checksum when available, local
manifest, class-selection rule, and sample limit. The lecture pipeline selects
the first fixed number of lexicographically sorted files in each class. This
makes membership deterministic across supported operating systems and balances
the classes, but the limit changes sampling uncertainty and therefore belongs
in every reported result.

IMDB supplies distinct \code{train} and \code{test} directories. We preserve
that boundary and create a validation split only from publisher-designated
training reviews:
\[
  D_{\mathrm{publisher\ train}}
  =D_{\mathrm{optimization}}\mathbin{\dot\cup}D_{\mathrm{validation}},
  \qquad
  D_{\mathrm{test}}\text{ remains untouched.}
\]
Stratification preserves class proportions, while a fixed split seed makes
membership reproducible. The optimization subset fits parameters and
preprocessing state. Validation data selects checkpoints, architectures, and
decision policy. Test data estimates the frozen procedure once.

\begin{professionalbox}
The automated test suite sets an explicit CI flag and replaces the download
with locally authored sentences. That fixture answers ``does the pipeline still
execute and obey its shape contracts?'' It cannot answer ``does this model
classify movie reviews?'' Outside CI, acquisition failure must raise an error
rather than silently substitute the fixture and change the scientific question.
\end{professionalbox}

\section{How text models developed}

Early automatic text classifiers commonly combined manually constructed
lexicons with counts of words or short word sequences. Statistical language
models and linear classifiers made these counts trainable and testable.
Distributed representations later replaced one coordinate per vocabulary word
with compact learned vectors. Convolutional networks detected local phrase
patterns, while recurrent networks propagated state across a sequence.
Attention mechanisms then allowed representations to combine information from
distant positions more directly, and large pretrained transformers made it
possible to adapt representations learned from vast corpora.

This history is not a sequence in which every new method makes earlier methods
obsolete. Count-based linear models remain valuable baselines. Small embedding
models can be appropriate when latency, energy, privacy, or data volume rule
out a large pretrained model. Every approach still depends on data provenance,
split integrity, metric choice, error analysis, and operational monitoring.

\section{The representation pipeline}

A raw string must cross several conceptually different boundaries before it
can enter a dense neural network:
\[
\text{raw review}
\longrightarrow \text{normalized text}
\longrightarrow \text{tokens}
\longrightarrow \text{integer IDs}
\longrightarrow \text{embeddings}
\longrightarrow \text{review vector}.
\]
Collapsing these steps under the word ``tokenization'' hides important modeling
decisions.

In TensorFlow, \code{TextVectorization} is a stateful preprocessing layer. Its
\code{adapt} method counts tokens and constructs a vocabulary; calling the
layer applies the frozen mapping. The distinction is analogous to
\code{fit} versus \code{transform} in a tabular pipeline. Only optimization
text may enter \code{adapt}. Validation, test, and future text call the already
adapted layer.

Keeping this layer inside the Keras model makes the deployed input contract a
batch of strings rather than a batch of unexplained integers. It also packages
normalization rules and vocabulary order with the fitted network. A parity test
should confirm that a saved and reloaded model returns the same logits for
versioned review fixtures.

\begin{misconceptionbox}{token IDs are not numerical measurements}
Assigning the word \emph{excellent} ID 200 and \emph{awful} ID 100 does not make
the first word twice the second. The IDs are addresses in a vocabulary. A
trainable embedding converts those addresses into vectors whose geometry is
learned from the training objective; treating raw IDs as scalar features would
invent an ordering and distance that language does not possess.
\end{misconceptionbox}

\subsection{Normalization}

Normalization may lowercase text, standardize Unicode, remove markup, or strip
punctuation. Each transformation encodes an assumption. Exclamation marks,
capitalization, emoji, spelling variants, and repeated characters can carry
sentiment. A default that is sensible for movie reviews may be destructive for
clinical notes, source code, or multilingual messages.

A normalization function should be deterministic, versioned, and identical in
training and inference. It should also be tested on empty strings, unusual
Unicode, markup, and adversarially long inputs.

\subsection{Tokens and vocabularies}

A tokenizer divides normalized text into units. Word tokenization is intuitive
but creates rare and unseen words. Character tokenization has a small
vocabulary but produces long sequences. Subword methods construct reusable
pieces, allowing an unfamiliar word to share structure with familiar ones.

Let the training vocabulary be
\[
  \mathcal V=(v_0,v_1,\ldots,v_{V-1}).
\]
An integerizer maps a token to its index in $\mathcal V$. These integers are
identifiers, not quantitative measurements: token 200 is not twice token 100.
A typical contract reserves index zero for padding and index one for an
out-of-vocabulary token, often displayed as \code{[UNK]}.

If $n_{\mathrm{OOV}}$ observed training tokens map to ID one and
$n_{\mathrm{observed}}$ excludes padding, then
\[
  r_{\mathrm{OOV}}=
  \frac{n_{\mathrm{OOV}}}{n_{\mathrm{observed}}}
\]
is a useful representation diagnostic. It is not a performance metric. A high
rate may indicate an undersized vocabulary, unexpected language, or domain
shift. Enlarging the vocabulary increases memory and can admit rare noisy
features, so the response should follow diagnosis rather than a fixed rule.

\begin{proposition}
If a vocabulary is selected using validation or test documents, then the
evaluation pipeline is not independent of those held-out documents.
\end{proposition}

\begin{proof}
Vocabulary selection is a learned transformation because the selected tokens
and their indices depend on observed token frequencies. If held-out documents
contribute to those frequencies, then the representation function supplied to
the estimator is a function of the held-out inputs. Consequently the fitted
pipeline has received information about the evaluation distribution before
evaluation, even if it has not received the held-out labels. The evaluation no
longer estimates a pipeline constructed from training data alone.
\end{proof}

This leakage may appear weak because labels are absent, but it breaks the
experimental contract. Split observations first. Fit normalization state,
vocabulary, and any other learned preprocessing using only the training
partition. Apply the frozen transformation to validation, test, and production
text.

\subsection{Padding, truncation, and masking}

Reviews have different lengths. A tensor batch is usually rectangular, so a
pipeline selects a sequence length $L$. Short sequences receive padding;
tokens beyond $L$ are truncated. These operations are not harmless bookkeeping.
A small $L$ may remove a final conclusion that reverses earlier language. A
large $L$ increases memory and computation.

A mask is a Boolean tensor marking positions that contain actual tokens rather
than padding. A compatible embedding or sequence layer uses the mask so that
padding does not contribute to a pooled representation or recurrent state.
Without masking, document length can become an unintended predictive feature.

Before choosing $L$, examine the training length distribution and state the
truncation rate. In deployment, monitor this rate because longer documents can
signal distribution shift or an input-contract violation.

For a batch of $B$ reviews, the representation shapes are
\[
\begin{array}{ccl}
\text{raw strings} &:& (B),\\
\text{token IDs} &:& (B,L),\\
\text{embedding vectors} &:& (B,L,d),\\
\text{padding mask} &:& (B,L),\\
\text{pooled review vectors} &:& (B,d),\\
\text{binary logits} &:& (B,1).
\end{array}
\]
Shape assertions make these interfaces executable. They catch accidental
scalar inputs, missing batch dimensions, wrong output width, and incompatible
label tensors before a long training run.

\section{Embeddings are trainable lookup tables}

Let $V$ denote vocabulary size and $d$ embedding dimension. An embedding layer
stores
\[
  E\in\mathbb R^{V\times d}.
\]
For token ID $i$, the layer returns row $E_{i,:}$. A length-$L$ sequence thus
becomes a matrix in $\mathbb R^{L\times d}$. The rows of $E$ begin from an
initializer and are updated by backpropagation because they participate in the
classification loss.

An embedding coordinate does not come with a fixed linguistic interpretation.
If two tokens acquire nearby vectors, they were useful in similar predictive
contexts under this corpus, objective, model, and training run. That geometry
does not establish causal, psychological, or universal semantic equivalence.

\subsection{Pooling a sequence}

The simplest fixed-width representation is a masked average. Let
$m_t=\mathbf 1[i_t\neq0]$ indicate a non-padding token. Define
\[
  h(x)=
  \frac{\sum_{t=1}^{L}m_tE_{i_t,:}}{\sum_{t=1}^{L}m_t}.
\]
This formula assumes at least one observed token. The input contract must
reject, abstain on, or deliberately represent an empty normalized input rather
than letting a zero denominator produce an undefined result.
For an admissible review it produces $h(x)\in\mathbb R^d$ regardless of review
length. A hidden layer
and output logit can then compute
\[
  r=\operatorname{ReLU}(W_hh(x)+b_h),
  \qquad
  z=w_o^\top r+b_o.
\]

Mean pooling is fast and makes tensor shapes transparent. It also largely
discards order. The phrases ``not painfully slow'' and ``not moving, painfully
slow'' can contain similar tokens but express different judgments. A
one-dimensional convolution can detect local patterns; a recurrent network can
propagate ordered state; and attention can form context-dependent combinations
across positions. The added structure should be justified with a controlled
comparison rather than assumed to help.

In Keras, \code{Embedding(mask\_zero=True)} constructs the Boolean mask and a
compatible \code{GlobalAveragePooling1D} consumes it. Padding ID zero is
ignored; OOV ID one is not. The latter represents observed but unresolved
language and should therefore contribute its learned OOV vector.

\section{A binary logit and a stable objective}

The model emits a logit $z\in\mathbb R$. The sigmoid converts this unrestricted
score to a value in $(0,1)$:
\[
  p_\theta(y=1\mid x)=\sigma(z)=\frac{1}{1+\exp(-z)}.
\]
For $y\in\{0,1\}$, binary cross-entropy is
\[
  \ell(z,y)
  =-y\log\sigma(z)-(1-y)\log(1-\sigma(z)).
\]
Mature frameworks should consume logits directly because a stable
implementation avoids first forming probabilities that round to exactly zero
or one.

\begin{proposition}
The derivative of binary cross-entropy with respect to its logit is
\[
  \frac{\partial\ell}{\partial z}=\sigma(z)-y.
\]
\end{proposition}

\begin{proof}
Let $p=\sigma(z)$. Since $dp/dz=p(1-p)$,
\[
\begin{aligned}
\frac{\partial\ell}{\partial z}
&=\left(-\frac{y}{p}+\frac{1-y}{1-p}\right)p(1-p)\\
&=-y(1-p)+(1-y)p\\
&=p-y.
\end{aligned}
\]
\end{proof}

The expression is the familiar prediction-minus-target error from logistic
regression. The difference is that $z=f_\theta(x)$ now depends on vectorization,
embedding lookup, pooling, and hidden layers. Automatic differentiation applies
the chain rule from $p-y$ through this graph.

The output and loss form one contract. A final dense layer with no activation
emits logits and must pair with
\code{BinaryCrossentropy(from\_logits=True)}. If the model instead emits
$\sigma(z)$, the loss must be configured for probabilities. Applying sigmoid
twice changes the objective and weakens gradients. Before fitting, evaluate one
batch and assert that the logit shape is $(B,1)$ and the loss is finite.

Keras's binary-accuracy threshold is expressed in the coordinate system of its
input. When the metric receives logits, threshold zero corresponds to
$\sigma(0)=0.5$. Confusing a logit threshold of zero with a probability
threshold of zero can make an apparently reasonable metric configuration
meaningless.

\begin{workedexample}{a TensorFlow sentiment pipeline}
A reproducible first model has the following explicit contract:
\begin{enumerate}
  \item preserve the publisher's train/test partition and reserve validation
  observations from the training partition;
  \item adapt a \code{TextVectorization} layer on training reviews only;
  \item output a fixed-length integer sequence with documented padding, OOV,
  vocabulary-size, and truncation policies;
  \item map IDs through an \code{Embedding} layer with masking enabled;
  \item use masked global-average pooling, a small ReLU hidden layer, and one
  output logit;
  \item optimize logits-based binary cross-entropy with Adam;
  \item select checkpoints and any threshold using validation data; and
  \item evaluate the frozen candidate once on test data.
\end{enumerate}

The baseline should record its seed, package versions, vocabulary, sequence
length, parameter count, epoch times, loss curves, and evaluation metrics.
The vectorizer belongs inside or beside the saved model artifact so that
training and serving cannot silently disagree about integer IDs.
\end{workedexample}

\section{Thresholds turn modeled probabilities into decisions}

The conventional classifier
\[
  \widehat y=\mathbf 1[\sigma(z)\geq0.5]
\]
is equivalently $\mathbf 1[z\geq0]$. Neither threshold is mandated by
cross-entropy. A threshold encodes an operational tradeoff.

For the positive class, precision and recall are
\[
  \operatorname{precision}=\frac{TP}{TP+FP},
  \qquad
  \operatorname{recall}=\frac{TP}{TP+FN}.
\]
If falsely promoting an unfavorable review is costly, precision may matter. If
missing critical feedback is costly, recall may matter. A confusion matrix
retains all four error counts; accuracy alone collapses them.

Threshold selection belongs on validation data and should reflect a stated
cost or capacity constraint. Repeatedly choosing a threshold after examining
test results converts the test set into another validation set. Probability
calibration should also be checked when consumers interpret scores as
probabilities rather than ranks.

A valid sequence is therefore:
\begin{enumerate}
  \item fit parameters on optimization data;
  \item select the checkpoint from validation loss;
  \item inspect validation precision--recall tradeoffs and freeze a threshold;
  \item compute test logits once, transform them with sigmoid, and apply that
  frozen threshold; and
  \item report the confusion counts, metrics, sample size, and uncertainty.
\end{enumerate}
For the classroom example no real operating cost is supplied, so the notebook
predeclares 0.5 instead of reverse-engineering a favorable threshold. A real
application must elicit costs, capacity, or a target constraint from its
decision context.

\conceptcheck{Suppose a team chooses a threshold after repeatedly inspecting
test-set recall. Which experimental boundary has been broken, and what data
should the team use to choose the threshold instead?}

\section{Error analysis is part of modeling}

Inspect false positives, false negatives, and low-margin correct predictions.
Useful error categories include:
\begin{itemize}
  \item negation and contrast, such as ``slow at first, but not disappointing'';
  \item sarcasm or rhetorical language;
  \item mixed evaluations of acting, direction, and plot;
  \item names, genres, or years used as shortcuts;
  \item influential tokens mapped to OOV;
  \item decisive evidence removed by truncation; and
  \item disagreement or ambiguity in the original rating rule.
\end{itemize}

An error category suggests an experiment. Negation may motivate order-sensitive
features. Truncation errors may motivate a larger $L$ or hierarchical model.
Genre shortcuts may require stratified evaluation. Ambiguous labels may place
an upper limit on meaningful performance. A larger architecture is not the
automatic response to every error.

\begin{professionalbox}
An error audit should retain the true label, frozen-model probability, decision,
error type, retained-token count, OOV count, and a bounded review excerpt when
the data policy permits inspection. Assign a hypothesized category before
changing the model. Then propose a controlled experiment and a result that
would falsify the hypothesis. Error browsing without a coding scheme easily
becomes anecdotal model tuning.
\end{professionalbox}

\section{Responsible use and deployment}

A movie-review classifier is a bounded educational task. Sentiment models are
often redeployed to employee messages, student writing, clinical narratives,
or customer support. That transfer changes the population, language, privacy
risk, and consequences. Performance on IMDB does not justify such use.

Potential failures include dialect-dependent error rates, changing slang,
spurious correlations with identity terms, forced binary predictions on mixed
sentiment, and exposure of personal text through logs. Evaluation should be
sliced by relevant, ethically supportable groups and linguistic phenomena,
with uncertainty reported for small samples. High-impact uses require human
review, appeal paths, access controls, retention limits, and possibly a decision
not to deploy.

A released text system should version together:
\begin{itemize}
  \item a data manifest and split identifiers;
  \item normalization and tokenization code;
  \item vocabulary order, OOV policy, and sequence length;
  \item architecture, parameters, optimizer state when resuming, and framework
  versions;
  \item class order, logit/probability contract, and decision threshold;
  \item validation and test evidence; and
  \item a model card describing intended use and known limitations.
\end{itemize}

Production monitoring should include request validity, empty-input frequency,
token length, truncation rate, OOV rate, score and decision distributions,
latency, failures, and delayed outcome metrics when legitimate labels arrive.
Raw user text should not appear in general logs. Aggregate diagnostics or
carefully controlled examples reduce privacy and security exposure.

\section{Exercises}
\begin{enumerate}
  \item Give separate definitions of normalization, tokenization,
  integerization, padding, masking, and embedding. For each, name one failure
  mode.
  \item Explain why adapting a vocabulary on unlabeled test text is still
  leakage. Construct a small example in which the selected vocabulary changes.
  \item For a vocabulary of size $8{,}000$, embedding dimension 48, sequence
  length 200, and batch size 32, compute the embedding parameter count and
  output-tensor shape.
  \item Derive the binary-cross-entropy logit gradient. Explain why using a
  sigmoid output with a loss configured for logits is incorrect.
  \item Compare masked mean pooling with a one-dimensional convolution while
  holding the split, vocabulary, sequence length, and evaluation protocol
  fixed. Report parameter count, runtime, validation loss, and error categories.
  \item Evaluate sequence lengths 100, 250, and 500. Plot truncation rate,
  runtime, and negative-class recall. Recommend a length using evidence rather
  than accuracy alone.
  \item Choose an operational setting for which false positives and false
  negatives have different costs. Define a threshold-selection rule without
  touching the test set.
  \item Conduct a ten-example error audit. Propose one data change and one model
  change, and state what result would falsify each proposed improvement.
  \item Write tests verifying reserved token IDs, output shape, finite logits,
  deterministic inference, empty-string handling, artifact reload, and parity
  between batch and single-review predictions.
  \item Draft a monitoring plan that detects changing review length, OOV rate,
  score distribution, and delayed accuracy without logging raw reviews.
\end{enumerate}

\begin{chapterreview}
Text classification begins with a construct and data contract, not an
architecture. Normalization, tokenization, vocabulary construction, padding,
and masking determine what evidence a model can use. Learned preprocessing
must fit the training partition alone. Embeddings replace arbitrary integer
identifiers with trainable vectors, while pooling or sequence layers produce a
fixed representation. A binary logit trained by cross-entropy retains the
logistic-regression error $\sigma(z)-y$; backpropagation distributes that error
through the text representation. Threshold metrics, direct error inspection,
artifact versioning, privacy-conscious monitoring, and responsible scope are
part of the model, not postscript engineering.
\end{chapterreview}

\begin{notebooklink}
Lecture 12 notebook 04 executes this full pipeline. The normal student route
downloads the genuine Stanford IMDB corpus from its publisher through Keras's
cache-aware data API. Only the automated CI route selects a bounded locally
authored fixture, and its results are explicitly excluded from substantive
analysis. Students inspect length and
label distributions, fit a training-only vocabulary, trace tensor shapes,
train an embedding classifier, vary thresholds, and audit held-out errors.
\end{notebooklink}
